\documentclass[12pt]{article}
\usepackage[margin=1in]{geometry}
\usepackage{setspace}
\onehalfspacing

\begin{document}
\thispagestyle{empty}

June 4, 2026

\vspace{0.5cm}

The Editors\\
\textit{Eurosurveillance}

\vspace{0.5cm}

Dear Editors,

\vspace{0.3cm}

We are pleased to submit our manuscript, \textbf{``Estimating Infectious Disease Importation Risk during the 2026 FIFA World Cup in the United States, June--July 2026,''} for consideration in \textit{Eurosurveillance}.

The 2026 FIFA World Cup --- the largest in the tournament's history --- will bring visitors from 48 nations to 11 US host cities between June 11 and July 19, 2026. We developed a Poisson importation framework applied to five diseases (dengue fever, influenza, malaria, measles, and pertussis), combining WHO surveillance data, BTS T-100 routing records, and CBP I-94 arrival volumes with Monte Carlo uncertainty propagation. Three nested travel models of increasing resolution allow us to isolate the contribution of World Cup fan travel from background tourism and to route risk to specific host cities based on match scheduling. All inputs are publicly available, making the analysis fully reproducible and extensible to future events.

Key findings include dengue as the dominant risk at most venue cities, driven by Brazil and Latin American nations; malaria exceeding dengue probability at Atlanta due to direct travel from West and Central Africa; and pertussis and measles concentrated among European source nations. Importantly, background tourism accounts for the dominant share of importation risk, suggesting that public health surveillance should not focus exclusively on fan travel streams.

We believe this work is well suited to \textit{Eurosurveillance} given the journal's commitment to timely, actionable infectious disease surveillance and its established record of publishing mass gathering risk assessments. Given the imminent start of the tournament on June 11, we respectfully request expedited editorial consideration.

This manuscript is original, has not been published previously, and is not under consideration at any other journal. All authors have approved the submission and declare no conflicts of interest. This study used only publicly available, de-identified data and did not require ethical approval.

\vspace{0.5cm}

Sincerely,

\vspace{0.3cm}

Jose L. Herrera-Diestra, PhD\\
Tarleton State University, Stephenville, TX, USA\\
jherreradiestra@tarleton.edu\\

\vspace{0.2cm}

\noindent On behalf of all co-authors:\\
Kaiming Bi (UTHealth Houston),
Susan Ptak (University of Texas at Austin),
Zeynep Ertem (Binghamton University),
Amera Al-amery (Jordan University of Science and Technology),
Mallory Harris (University of Texas at Austin),
Lauren Ancel Meyers (University of Texas at Austin)

\end{document}
